\textbf{Phần 1. Trắc nghiệm lựa chọn}
\setcounterpageref{ex}{0}

\begin{ex} %[5D12V72]
Một xe tải chở $2{,}5\text{ tấn}$ gạo đi giao cho các đại lý. Bác tài đã giao ở cửa hàng thứ nhất $850\text{ kg}$ và cửa hàng thứ hai $950\text{ kg}$. Số gạo còn lại trên xe tải là bao nhiêu tạ?
\choice
{$6\text{ tạ}$}
{\True $7\text{ tạ}$}
{$8\text{ tạ}$}
{$70\text{ tạ}$}
\loigiai{
Đổi: $2{,}5\text{ tấn} = 2\,500\text{ kg}$.\\
Số ki-lô-gam gạo đã giao ở hai cửa hàng là:
$$850 + 950 = 1\,800\text{ (kg)}$$
Số gạo còn lại trên xe là:
$$2\,500 - 1\,800 = 700\text{ (kg)} = 7\text{ (tạ)}$$
\textbf{Đáp án của câu hỏi là B}
}
\end{ex}

\begin{ex} %[5D12V72]
Một thửa ruộng có diện tích $0{,}4\text{ ha}$. Trung bình cứ mỗi $100\text{ m}^2$ thu hoạch được $65\text{ kg}$ thóc. Hỏi cả thửa ruộng đó thu hoạch được bao nhiêu tấn thóc?
\choice
{$26\text{ tấn}$}
{\True $2{,}6\text{ tấn}$}
{$260\text{ tấn}$}
{$0{,}26\text{ tấn}$}
\loigiai{
Đổi: $0{,}4\text{ ha} = 4\,000\text{ m}^2$.\\
Số ki-lô-gam thóc thu hoạch được là:
$$4\,000 : 100 \times 65 = 2\,600\text{ (kg)} = 2{,}6\text{ (tấn)}$$
\textbf{Đáp án của câu hỏi là B}
}
\end{ex}

\begin{ex} %[5D12V72]
Một thợ may làm việc may xong $6$ chiếc áo sơ mi hết $4$ giờ $30$ phút. Biết thời gian may mỗi chiếc áo là như nhau. Hỏi người đó may mỗi chiếc áo hết bao nhiêu phút?
\choice
{$40\text{ phút}$}
{\True $45\text{ phút}$}
{$50\text{ phút}$}
{$35\text{ phút}$}
\loigiai{
Đổi: $4\text{ giờ } 30\text{ phút} = 270\text{ phút}$.\\
Thời gian may mỗi chiếc áo là:
$$270 : 6 = 45\text{ (phút)}$$
\textbf{Đáp án của câu hỏi là B}
}
\end{ex}

\begin{ex} %[5D12V72]
Một bể nước dự trữ có dung tích là $3{,}6\text{ m}^3$. Mỗi ngày một hộ gia đình sử dụng hết trung bình $450\text{ l}$ nước. Hỏi khi bể đầy nước thì đủ cho gia đình đó sử dụng trong bao nhiêu ngày?
\choice
{$6\text{ ngày}$}
{\True $8\text{ ngày}$}
{$10\text{ ngày}$}
{$12\text{ ngày}$}
\loigiai{
Đổi: $3{,}6\text{ m}^3 = 3\,600\text{ dm}^3 = 3\,600\text{ l}$.\\
Số ngày đủ để gia đình sử dụng là:
$$3\,600 : 450 = 8\text{ (ngày)}$$
\textbf{Đáp án của câu hỏi là B}
}
\end{ex}

\begin{ex} %[5D12V72]
Một sợi dây đồng dài $12\text{ m } 8\text{ dm}$. Người ta cắt sợi dây đó thành các đoạn ngắn bằng nhau, mỗi đoạn dài $16\text{ cm}$. Số đoạn dây cắt được là
\choice
{$60\text{ đoạn}$}
{\True $80\text{ đoạn}$}
{$90\text{ đoạn}$}
{$75\text{ đoạn}$}
\loigiai{
Đổi: $12\text{ m } 8\text{ dm} = 1\,280\text{ cm}$.\\
Số đoạn dây cắt được là:
$$1\,280 : 16 = 80\text{ (đoạn)}$$
\textbf{Đáp án của câu hỏi là B}
}
\end{ex}

\begin{ex} %[5D12V72]
Một thùng chứa $50$ chai dầu ăn, mỗi chai chứa $0{,}75\text{ l}$ dầu. Biết mỗi lít dầu ăn nặng $0{,}8\text{ kg}$ và vỏ thùng rỗng nặng $2\text{ kg}$. Toàn bộ thùng dầu ăn đó nặng bao nhiêu ki-lô-gam?
\choice
{$30\text{ kg}$}
{\True $32\text{ kg}$}
{$34\text{ kg}$}
{$40\text{ kg}$}
\loigiai{
Tổng số lít dầu ăn trong thùng là:
$$50 \times 0{,}75 = 37{,}5\text{ (l)}$$
Khối lượng dầu ăn trong thùng là:
$$37{,}5 \times 0{,}8 = 30\text{ (kg)}$$
Khối lượng toàn bộ thùng dầu là:
$$30 + 2 = 32\text{ (kg)}$$
\textbf{Đáp án của câu hỏi là B}
}
\end{ex}

\begin{ex} %[5D12V72]
Một sân chơi hình chữ nhật có diện tích $180\text{ m}^2$. Người ta lát sân bằng các viên gạch men hình vuông có cạnh $60\text{ cm}$. Số viên gạch men cần dùng để lát kín sân chơi đó là
\choice
{$400\text{ viên}$}
{\True $500\text{ viên}$}
{$600\text{ viên}$}
{$300\text{ viên}$}
\loigiai{
Đổi: $60\text{ cm} = 0{,}6\text{ m}$.\\
Diện tích một viên gạch men là:
$$0{,}6 \times 0{,}6 = 0{,}36\text{ (m}^2\text{)}$$
Số viên gạch men cần dùng là:
$$180 : 0{,}36 = 500\text{ (viên)}$$
\textbf{Đáp án của câu hỏi là B}
}
\end{ex}

\begin{ex} %[5D12V72]
Cho bốn số đo khối lượng sau: $A = 3\text{ tấn } 45\text{ kg}$; $B = 3{,}4\text{ tấn}$; $C = 30{,}5\text{ tạ}$; $D = 3\,050\text{ kg}$. Số đo lớn nhất trong bốn số đo trên là
\choice
{$A$}
{\True $B$}
{$C$}
{$D$}
\loigiai{
Đổi các số đo về cùng đơn vị ki-lô-gam:
\begin{itemize}
    \item $A = 3\text{ tấn } 45\text{ kg} = 3\,045\text{ kg}$.
    \item $B = 3{,}4\text{ tấn} = 3\,400\text{ kg}$.
    \item $C = 30{,}5\text{ tạ} = 3\,050\text{ kg}$.
    \item $D = 3\,050\text{ kg}$.
\end{itemize}
Vì $3\,400\text{ kg}$ lớn nhất nên số đo lớn nhất là $B = 3{,}4\text{ tấn}$.\\
\textbf{Đáp án của câu hỏi là B}
}
\end{ex}

\begin{ex} %[5D12V72]
Một vòi nước chảy vào một bể có dung tích $4\text{ m}^3$. Mỗi phút vòi chảy được $30\text{ l}$ nước. Sau $1$ giờ $40$ phút chảy liên tục, lượng nước đã chảy vào bể chiếm bao nhiêu phần trăm dung tích của bể?
\choice
{$70\%$}
{\True $75\%$}
{$80\%$}
{$85\%$}
\loigiai{
Đổi: $4\text{ m}^3 = 4\,000\text{ l}$; $1\text{ giờ } 40\text{ phút} = 100\text{ phút}$.\\
Lượng nước đã chảy vào bể sau $100$ phút là:
$$30 \times 100 = 3\,000\text{ (l)}$$
Tỉ số phần trăm lượng nước so với dung tích bể là:
$$3\,000 : 4\,000 = 0{,}75 = 75\%$$
\textbf{Đáp án của câu hỏi là B}
}
\end{ex}

\begin{ex} %[5D12V72]
Một chuyến tàu hỏa xuất phát từ Hà Nội lúc $21$ giờ $15$ phút đêm hôm trước và đến thành phố Huế lúc $10$ giờ $45$ phút trưa hôm sau. Thời gian chuyến tàu chạy trên đường là
\choice
{$12\text{ giờ } 30\text{ phút}$}
{\True $13\text{ giờ } 30\text{ phút}$}
{$14\text{ giờ } 15\text{ phút}$}
{$13\text{ giờ } 45\text{ phút}$}
\loigiai{
Thời gian tàu chạy từ $21$ giờ $15$ phút đến $24$ giờ là:
$$24\text{ giờ} - 21\text{ giờ } 15\text{ phút} = 2\text{ giờ } 45\text{ phút}$$
Tổng thời gian tàu chạy là:
$$2\text{ giờ } 45\text{ phút} + 10\text{ giờ } 45\text{ phút} = 13\text{ giờ } 30\text{ phút}$$
\textbf{Đáp án của câu hỏi là B}
}
\end{ex}

\begin{ex} %[5D12V72]
Một bể nước có đáy hình vuông cạnh $1{,}2\text{ m}$. Người ta đổ vào bể $720\text{ l}$ nước. Mực nước trong bể dâng cao thêm bao nhiêu đề-xi-mét?
\choice
{$4\text{ dm}$}
{\True $5\text{ dm}$}
{$6\text{ dm}$}
{$4{,}5\text{ dm}$}
\loigiai{
Đổi: $1{,}2\text{ m} = 12\text{ dm}$; $720\text{ l} = 720\text{ dm}^3$.\\
Diện tích đáy bể là:
$$12 \times 12 = 144\text{ (dm}^2\text{)}$$
Mực nước dâng cao thêm là:
$$720 : 144 = 5\text{ (dm)}$$
\textbf{Đáp án của câu hỏi là B}
}
\end{ex}

\begin{ex} %[5D12V72]
Đoạn đường thứ nhất dài $2\text{ km } 400\text{ m}$. Đoạn đường thứ hai dài bằng $75\%$ chiều dài đoạn đường thứ nhất. Tổng chiều dài của cả hai đoạn đường là bao nhiêu ki-lô-mét?
\choice
{$3{,}6\text{ km}$}
{\True $4{,}2\text{ km}$}
{$4{,}5\text{ km}$}
{$4\text{ km}$}
\loigiai{
Đổi: $2\text{ km } 400\text{ m} = 2{,}4\text{ km}$.\\
Chiều dài đoạn đường thứ hai là:
$$2{,}4 \times 75 : 100 = 1{,}8\text{ (km)}$$
Tổng chiều dài cả hai đoạn đường là:
$$2{,}4 + 1{,}8 = 4{,}2\text{ (km)}$$
\textbf{Đáp án của câu hỏi là B}
}
\end{ex}

\textbf{Phần 2. Trắc nghiệm đúng sai}
\setcounterpageref{ex}{0}

\begin{ex} %[5D12V72]
Một thửa ruộng hình chữ nhật có chiều dài $80\text{ m}$, chiều rộng $50\text{ m}$. Xét tính đúng sai của các khẳng định sau:
\choiceTF
{\True Diện tích của thửa ruộng là $4\,000\text{ m}^2$, tương đương $0{,}4\text{ ha}$.}
{\True Nếu cứ $100\text{ m}^2$ thu hoạch được $70\text{ kg}$ thóc thì cả thửa ruộng thu hoạch được $2{,}8\text{ tấn}$ thóc.}
{Nếu bán toàn bộ số thóc đó với giá $9\,000\text{ đồng/kg}$ thì thu được $28\text{ triệu đồng}$.}
{\True Số tiền bán thóc thu được là $25\,200\,000\text{ đồng}$.}
\loigiai{
\begin{itemchoice}
    \itemch Đúng. Diện tích: $80 \times 50 = 4\,000\text{ m}^2 = 0{,}4\text{ ha}$.
    \itemch Đúng. Khối lượng thóc: $4\,000 : 100 \times 70 = 2\,800\text{ kg} = 2{,}8\text{ tấn}$.
    \itemch Sai. Số tiền thu được: $2\,800 \times 9\,000 = 25\,200\,000\text{ đồng} \neq 28\,000\,000\text{ đồng}$.
    \itemch Đúng. Số tiền thu được là $25\,200\,000\text{ đồng}$.
\end{itemchoice}
}
\end{ex}

\begin{ex} %[5D12V72]
Một vận động viên hoàn thành $4$ vòng chạy tiếp sức trên sân vận động với tổng thời gian là $6$ phút $20$ giây. Xét tính đúng sai của các khẳng định sau:
\choiceTF
{\True Tổng thời gian chạy của vận động viên đổi ra giây là $380\text{ giây}$.}
{\True Thời gian chạy trung bình mỗi vòng là $1$ phút $35$ giây.}
{Thời gian chạy trung bình mỗi vòng là $90\text{ giây}$.}
{\True Nếu duy trì tốc độ đó thì thời gian chạy $6$ vòng là $9$ phút $30$ giây.}
\loigiai{
\begin{itemchoice}
    \itemch Đúng. $6\text{ phút } 20\text{ giây} = 6 \times 60 + 20 = 380\text{ giây}$.
    \itemch Đúng. Thời gian trung bình 1 vòng: $380 : 4 = 95\text{ giây} = 1\text{ phút } 35\text{ giây}$.
    \itemch Sai. $95\text{ giây} \neq 90\text{ giây}$.
    \itemch Đúng. Thời gian chạy 6 vòng: $95 \times 6 = 570\text{ giây} = 9\text{ phút } 30\text{ giây}$.
\end{itemchoice}
}
\end{ex}

\textbf{Phần 3. Trả lời ngắn}
\setcounterpageref{ex}{0}

\begin{ex} %[5D12V72]
Đổi $2{,}45\text{ m}^3$ ra đơn vị đề-xi-mét khối.
\shortans{2450}
\loigiai{
Vì $1\text{ m}^3 = 1\,000\text{ dm}^3$ nên:
$$2{,}45 \times 1\,000 = 2\,450\text{ (dm}^3\text{)}$$
\textbf{Đáp số:} $2450$
}
\end{ex}

\begin{ex} %[5D12V72]
Một bao xi măng nặng $50\text{ kg}$. Cần bao nhiêu bao xi măng như thế để có tổng khối lượng là $3\text{ tấn}$?
\shortans{60}
\loigiai{
Đổi: $3\text{ tấn} = 3\,000\text{ kg}$.\\
Số bao xi măng cần dùng là:
$$3\,000 : 50 = 60\text{ (bao)}$$
\textbf{Đáp số:} $60$
}
\end{ex}

\begin{ex} %[5D12V72]
Đổi $3$ giờ $45$ phút ra đơn vị phút.
\shortans{225}
\loigiai{
Ta có:
$$3 \times 60 + 45 = 225\text{ (phút)}$$
\textbf{Đáp số:} $225$
}
\end{ex}

\begin{ex} %[5D12V72]
Một khu đất có diện tích $1{,}2\text{ ha}$. Người ta chia đều khu đất thành $15$ mảnh nhỏ bằng nhau. Tính diện tích mỗi mảnh đất theo mét vuông.
\shortans{800}
\loigiai{
Đổi: $1{,}2\text{ ha} = 12\,000\text{ m}^2$.\\
Diện tích mỗi mảnh đất nhỏ là:
$$12\,000 : 15 = 800\text{ (m}^2\text{)}$$
\textbf{Đáp số:} $800$
}
\end{ex}

\begin{ex} %[5D12V72]
Một xe tải chở $24$ kiện hàng, mỗi kiện hàng nặng $125\text{ kg}$. Tính tổng khối lượng hàng trên xe tải đó theo tạ.
\shortans{30}
\loigiai{
Tổng khối lượng hàng là:
$$24 \times 125 = 3\,000\text{ (kg)} = 30\text{ (tạ)}$$
\textbf{Đáp số:} $30$
}
\end{ex}

\begin{ex} %[5D12V72]
Một vòi nước chảy vào bể với lưu lượng $20\text{ l}$ mỗi phút. Hỏi sau bao nhiêu phút thì vòi chảy được lượng nước là $1{,}8\text{ m}^3$?
\shortans{90}
\loigiai{
Đổi: $1{,}8\text{ m}^3 = 1\,800\text{ l}$.\\
Thời gian cần để chảy được lượng nước đó là:
$$1\,800 : 20 = 90\text{ (phút)}$$
\textbf{Đáp số:} $90$
}
\end{ex}
